\chapter{Quantenzahlen und Erhaltungssaetze}

\section{Ziel dieses Kapitels}

Dieses Kapitel behandelt Quantenzahlen von Teilchen und die zugehoerigen
Erhaltungssaetze.

Quantenzahlen sind wichtig, weil sie entscheiden, ob eine Teilchenreaktion oder ein
Zerfall erlaubt ist. Ausserdem helfen sie dabei, die beteiligte Wechselwirkung zu
erkennen.

\begin{notebox}
Dieses Kapitel behandelt die Stichwoerter:
\begin{itemize}
    \item elektrische Ladung,
    \item Baryonenzahl,
    \item Leptonenzahl,
    \item Flavor-Quantenzahlen,
    \item Strangeness,
    \item Charm,
    \item Bottomness,
    \item Topness,
    \item Isospin,
    \item dritte Isospinkomponente,
    \item Hyperladung,
    \item Paritaet,
    \item Erhaltungssaetze,
    \item starke Wechselwirkung,
    \item elektromagnetische Wechselwirkung,
    \item schwache Wechselwirkung.
\end{itemize}
\end{notebox}

\section{Was sind Quantenzahlen?}

\begin{definitionbox}
\emph{Quantenzahlen} sind Zahlen, die bestimmte Eigenschaften eines Teilchens oder
Zustands beschreiben.
\end{definitionbox}

\begin{conceptbox}
Wichtige Quantenzahlen in der Teilchenphysik sind:
\begin{itemize}
    \item elektrische Ladung $Q$,
    \item Baryonenzahl $B$,
    \item Leptonenzahl $L$,
    \item Flavor-Quantenzahlen wie Strangeness $S$,
    \item Isospin $I$ und dritte Komponente $I_3$,
    \item Hyperladung $Y$,
    \item Spin $J$,
    \item Paritaet $P$.
\end{itemize}
\end{conceptbox}

\begin{conceptbox}
Eine Quantenzahl kann:
\begin{itemize}
    \item immer erhalten sein,
    \item nur bei bestimmten Wechselwirkungen erhalten sein,
    \item bei bestimmten Wechselwirkungen verletzt werden.
\end{itemize}

Genau deshalb sind Quantenzahlen ein Werkzeug, um Reaktionen zu analysieren.
\end{conceptbox}

\begin{examtipbox}
Bei einer Teilchenreaktion zuerst pruefen:
\begin{itemize}
    \item Ist elektrische Ladung erhalten?
    \item Ist Baryonenzahl erhalten?
    \item Ist Leptonenzahl erhalten?
    \item Sind Flavor-Quantenzahlen erhalten?
    \item Welche Wechselwirkung kann die beobachtete Aenderung erklaeren?
\end{itemize}
\end{examtipbox}

\section{Additive und multiplikative Quantenzahlen}

\begin{definitionbox}
Eine \emph{additive Quantenzahl} wird fuer ein Mehrteilchensystem durch Addition
der Einzelwerte berechnet.
\end{definitionbox}

\begin{examplebox}
Die elektrische Ladung ist additiv.

Fuer ein Proton mit Quarkinhalt
\[
    p=uud
\]
gilt:
\[
    Q_p
    =
    Q_u+Q_u+Q_d
    =
    \frac{2}{3}
    +
    \frac{2}{3}
    -
    \frac{1}{3}
    =
    1.
\]
\end{examplebox}

\begin{definitionbox}
Eine \emph{multiplikative Quantenzahl} wird fuer ein Mehrteilchensystem durch
Multiplikation der Einzelwerte berechnet.
\end{definitionbox}

\begin{conceptbox}
Die Paritaet ist ein Beispiel fuer eine multiplikative Quantenzahl.
Bei zusammengesetzten Systemen muss man neben den intrinsischen Paritaeten auch den
Bahndrehimpuls beruecksichtigen.
\end{conceptbox}

\section{Elektrische Ladung}

\begin{definitionbox}
Die \emph{elektrische Ladung} $Q$ gibt an, wie stark ein Teilchen an der
elektromagnetischen Wechselwirkung teilnimmt.
\end{definitionbox}

\begin{conceptbox}
Elektrische Ladung ist in allen bekannten fundamentalen Prozessen erhalten.
Das gilt fuer starke, elektromagnetische und schwache Wechselwirkung.
\end{conceptbox}

\begin{notebox}
Typische Ladungen:
\[
    Q(e^-)=-1,
    \qquad
    Q(e^+)=+1,
    \qquad
    Q(\nu)=0.
\]
\[
    Q(u)=+\frac{2}{3},
    \qquad
    Q(d)=-\frac{1}{3}.
\]
\end{notebox}

\begin{examtipbox}
Elektrische Ladung ist der erste und einfachste Check.
Wenn die Ladung nicht erhalten ist, ist die Reaktion so nicht erlaubt.
\end{examtipbox}

\section{Baryonenzahl}

\begin{definitionbox}
Die \emph{Baryonenzahl} $B$ ist eine additive Quantenzahl.

Quarks haben:
\[
    B=\frac{1}{3}.
\]

Antiquarks haben:
\[
    B=-\frac{1}{3}.
\]
\end{definitionbox}

\begin{conceptbox}
Baryonen bestehen im einfachen Quarkmodell aus drei Quarks.
Daher gilt fuer Baryonen:
\[
    B=1.
\]

Antibaryonen bestehen aus drei Antiquarks.
Daher gilt fuer Antibaryonen:
\[
    B=-1.
\]

Mesonen bestehen aus einem Quark und einem Antiquark.
Daher gilt fuer Mesonen:
\[
    B=0.
\]
\end{conceptbox}

\begin{examplebox}
Proton:
\[
    p=uud.
\]
Dann:
\[
    B_p
    =
    \frac{1}{3}
    +
    \frac{1}{3}
    +
    \frac{1}{3}
    =
    1.
\]

Pion:
\[
    \pi^+ = u\overline{d}.
\]
Dann:
\[
    B_{\pi^+}
    =
    \frac{1}{3}
    -
    \frac{1}{3}
    =
    0.
\]
\end{examplebox}

\begin{examtipbox}
Baryonenzahl ist in normalen Teilchenreaktionen erhalten.

Ein Proton kann daher nicht einfach nur in ein Positron und Photonen zerfallen,
weil dabei die Baryonenzahl verschwinden wuerde.
\end{examtipbox}

\section{Leptonenzahl}

\begin{definitionbox}
Die \emph{Leptonenzahl} $L$ ist eine additive Quantenzahl.

Leptonen haben:
\[
    L=+1.
\]

Antileptonen haben:
\[
    L=-1.
\]

Nichtleptonen haben:
\[
    L=0.
\]
\end{definitionbox}

\begin{conceptbox}
Man kann auch getrennte Familien-Leptonenzahlen definieren:
\[
    L_e,
    \qquad
    L_\mu,
    \qquad
    L_\tau.
\]

Dabei zaehlt $L_e$ Elektronen und Elektron-Neutrinos,
$L_\mu$ Myonen und Myon-Neutrinos,
und $L_\tau$ Tauonen und Tau-Neutrinos.
\end{conceptbox}

\begin{examplebox}
Beim Beta-Minus-Zerfall:
\[
    n \rightarrow p + e^- + \overline{\nu}_e.
\]

Leptonenzahl links:
\[
    L=0.
\]

Leptonenzahl rechts:
\[
    L(e^-)+L(\overline{\nu}_e)
    =
    1-1
    =
    0.
\]

Die Leptonenzahl ist erhalten.
\end{examplebox}

\begin{examtipbox}
Beim Beta-Zerfall ist die Leptonenzahl der Grund, warum beim Elektron ein
Antineutrino und beim Positron ein Neutrino auftritt.
\end{examtipbox}

\section{Flavor und Flavor-Quantenzahlen}

\begin{definitionbox}
Der \emph{Flavor} eines Quarks bezeichnet seine Sorte:
\[
    u,\quad d,\quad s,\quad c,\quad b,\quad t.
\]
\end{definitionbox}

\begin{conceptbox}
Zu bestimmten Quark-Flavors gehoeren Flavor-Quantenzahlen:
\begin{itemize}
    \item Strangeness $S$ fuer strange-Quarks,
    \item Charm $C$ fuer charm-Quarks,
    \item Bottomness beziehungsweise Beauty $B'$ fuer bottom-Quarks,
    \item Topness $T$ fuer top-Quarks.
\end{itemize}
\end{conceptbox}

\begin{notebox}
Achtung:
Die Baryonenzahl wird meistens mit $B$ bezeichnet.
Bottomness wird deshalb manchmal anders bezeichnet, zum Beispiel mit $B'$ oder
einfach verbal als Bottomness.
\end{notebox}

\section{Strangeness}

Strangeness ist historisch besonders wichtig, weil seltsame Teilchen stark erzeugt
werden koennen, aber schwach zerfallen.

\begin{definitionbox}
Die \emph{Strangeness} $S$ ist eine Flavor-Quantenzahl.

Fuer das strange-Quark gilt:
\[
    S(s)=-1.
\]

Fuer das strange-Antiquark gilt:
\[
    S(\overline{s})=+1.
\]
\end{definitionbox}

\begin{conceptbox}
Teilchen ohne strange-Quark und ohne strange-Antiquark haben:
\[
    S=0.
\]

Ein Teilchen mit einem strange-Quark hat typischerweise negative Strangeness.
Ein Teilchen mit einem strange-Antiquark hat positive Strangeness.
\end{conceptbox}

\begin{examplebox}
Das Kaon
\[
    K^+ = u\overline{s}
\]
hat:
\[
    S(K^+)=+1.
\]

Das Kaon
\[
    K^- = s\overline{u}
\]
hat:
\[
    S(K^-)=-1.
\]
\end{examplebox}

\begin{examplebox}
Das Lambda-Baryon hat Quarkinhalt:
\[
    \Lambda^0 = uds.
\]
Da ein strange-Quark enthalten ist, gilt:
\[
    S(\Lambda^0)=-1.
\]
\end{examplebox}

\begin{examtipbox}
Merke:
\[
    S(s)=-1,
    \qquad
    S(\overline{s})=+1.
\]
Das Vorzeichen ist eine haeufige Fehlerquelle.
\end{examtipbox}

\section{Charm, Bottomness und Topness}

\begin{definitionbox}
Fuer schwere Quarks definiert man analog weitere Flavor-Quantenzahlen:
\[
    C(c)=+1,
    \qquad
    C(\overline{c})=-1.
\]
\[
    B'(b)=-1,
    \qquad
    B'(\overline{b})=+1.
\]
\[
    T(t)=+1,
    \qquad
    T(\overline{t})=-1.
\]
\end{definitionbox}

\begin{notebox}
Bei Bottomness gibt es unterschiedliche Vorzeichenkonventionen in der Literatur.
Wichtig ist, in einer Aufgabe die verwendete Konvention konsistent zu benutzen.
\end{notebox}

\begin{conceptbox}
In der Vorlesung sind fuer viele Aufgaben vor allem Strangeness und die leichten
Hadronen wichtig.

Charmonium und Bottomonium werden spaeter bei Quarkonia behandelt.
Dort spielen charm- und bottom-Quarks eine zentrale Rolle.
\end{conceptbox}

\section{Erhaltung von Flavor}

\begin{conceptbox}
Starke und elektromagnetische Wechselwirkungen aendern den Quark-Flavor nicht.

Die schwache Wechselwirkung kann Quark-Flavors aendern.
Zum Beispiel:
\[
    d \rightarrow u + W^-.
\]
\end{conceptbox}

\begin{notebox}
Daraus folgt:
Flavor-Quantenzahlen wie Strangeness sind bei starker und elektromagnetischer
Wechselwirkung erhalten, koennen aber bei schwacher Wechselwirkung verletzt werden.
\end{notebox}

\begin{examplebox}
Starke Erzeugung seltsamer Teilchen erfolgt typischerweise paarweise:
\[
    \pi^- + p \rightarrow K^0 + \Lambda^0.
\]

Links:
\[
    S=0.
\]

Rechts:
\[
    S(K^0)+S(\Lambda^0)
    =
    +1 + (-1)
    =
    0.
\]

Die gesamte Strangeness bleibt erhalten.
\end{examplebox}

\begin{examplebox}
Ein schwacher Zerfall kann Strangeness aendern:
\[
    \Lambda^0 \rightarrow p + \pi^-.
\]

Links:
\[
    S(\Lambda^0)=-1.
\]

Rechts:
\[
    S(p)+S(\pi^-)=0+0=0.
\]

Also:
\[
    \Delta S = +1.
\]

Das ist ein Hinweis auf schwache Wechselwirkung.
\end{examplebox}

\begin{examtipbox}
Wenn Strangeness in einer Reaktion nicht erhalten ist, ist die Reaktion nicht stark
oder elektromagnetisch.

Dann ist sie ein Kandidat fuer die schwache Wechselwirkung.
\end{examtipbox}

\section{Isospin}

Isospin ist eine Quantenzahl, die Proton und Neutron beziehungsweise Up- und
Down-Quark in eine gemeinsame Struktur bringt.

\begin{definitionbox}
Der \emph{Isospin} $I$ ist eine Quantenzahl, die Teilchen mit aehnlichen Eigenschaften
unter der starken Wechselwirkung zu Multipletts zusammenfasst.
\end{definitionbox}

\begin{conceptbox}
Proton und Neutron bilden ein Isospin-Dublett:
\[
    N=
    \begin{pmatrix}
        p \\
        n
    \end{pmatrix}.
\]

Beide haben:
\[
    I=\frac{1}{2}.
\]

Die dritte Isospinkomponente unterscheidet sie:
\[
    I_3(p)=+\frac{1}{2},
    \qquad
    I_3(n)=-\frac{1}{2}.
\]
\end{conceptbox}

\begin{conceptbox}
Auch Up- und Down-Quark bilden ein Isospin-Dublett:
\[
    \begin{pmatrix}
        u \\
        d
    \end{pmatrix}.
\]

Dabei gilt:
\[
    I_3(u)=+\frac{1}{2},
    \qquad
    I_3(d)=-\frac{1}{2}.
\]
\end{conceptbox}

\begin{examtipbox}
Isospin ist besonders nuetzlich fuer die starke Wechselwirkung.

Die starke Wechselwirkung behandelt $u$ und $d$ naeherungsweise gleich.
Daher kann man sie als zwei Zustaende eines Isospin-Dubletts auffassen.
\end{examtipbox}

\section{Isospin-Multipletts}

\begin{definitionbox}
Ein \emph{Isospin-Multiplett} ist eine Gruppe von Teilchen mit gleichem Isospin $I$,
aber unterschiedlichen Werten von $I_3$.
\end{definitionbox}

\begin{conceptbox}
Fuer gegebenes $I$ gibt es:
\[
    2I+1
\]
verschiedene Werte von $I_3$.

Diese Werte laufen von:
\[
    -I
\]
bis
\[
    +I.
\]
\end{conceptbox}

\begin{examplebox}
Die Pionen bilden ein Isospin-Triplett:
\[
    \pi^+,\qquad \pi^0,\qquad \pi^-.
\]

Hier gilt:
\[
    I=1.
\]

Die drei Werte von $I_3$ sind:
\[
    +1,\qquad 0,\qquad -1.
\]
\end{examplebox}

\begin{examplebox}
Proton und Neutron bilden ein Isospin-Dublett:
\[
    p,\qquad n.
\]

Hier gilt:
\[
    I=\frac{1}{2}.
\]

Die zwei Werte von $I_3$ sind:
\[
    +\frac{1}{2},
    \qquad
    -\frac{1}{2}.
\]
\end{examplebox}

\section{Hyperladung}

Die Hyperladung verbindet Baryonenzahl und Strangeness.

\begin{definitionbox}
Die \emph{Hyperladung} $Y$ ist definiert als:
\[
    Y = B + S
\]
fuer Hadronen mit Strangeness.
\end{definitionbox}

\begin{notebox}
Allgemeiner kann man weitere Flavor-Quantenzahlen einbeziehen.
Im einfachen Kontext mit strange-Hadronen ist
\[
    Y=B+S
\]
die wichtigste Form.
\end{notebox}

\begin{conceptbox}
Die Hyperladung ist besonders wichtig fuer die Darstellung von Hadronen in
Multiplettdiagrammen.

Dort werden Teilchen haeufig in einem Diagramm mit
\[
    I_3
\]
auf der horizontalen Achse und
\[
    Y
\]
auf der vertikalen Achse angeordnet.
\end{conceptbox}

\section{Gell-Mann--Nishijima-Formel}

Elektrische Ladung, Isospin und Hyperladung haengen zusammen.

\begin{formulabox}
Die Gell-Mann--Nishijima-Formel lautet:
\[
    Q = I_3 + \frac{Y}{2}.
\]
\end{formulabox}

\begin{examplebox}
Fuer das Proton gilt:
\[
    I_3(p)=+\frac{1}{2},
    \qquad
    B=1,
    \qquad
    S=0.
\]
Damit:
\[
    Y=B+S=1.
\]
Also:
\[
    Q
    =
    I_3+\frac{Y}{2}
    =
    \frac{1}{2}
    +
    \frac{1}{2}
    =
    1.
\]
\end{examplebox}

\begin{examplebox}
Fuer das Neutron gilt:
\[
    I_3(n)=-\frac{1}{2},
    \qquad
    B=1,
    \qquad
    S=0.
\]
Damit:
\[
    Y=1.
\]
Also:
\[
    Q
    =
    -\frac{1}{2}
    +
    \frac{1}{2}
    =
    0.
\]
\end{examplebox}

\begin{examplebox}
Fuer das Lambda-Baryon gilt:
\[
    \Lambda^0=uds.
\]
Damit:
\[
    B=1,
    \qquad
    S=-1.
\]
Also:
\[
    Y=B+S=0.
\]
Da
\[
    I_3=0
\]
gilt, folgt:
\[
    Q=I_3+\frac{Y}{2}=0.
\]
\end{examplebox}

\begin{examtipbox}
Die Formel
\[
    Q=I_3+\frac{Y}{2}
\]
ist sehr nuetzlich, um Multipletts zu kontrollieren.
\end{examtipbox}

\section{Spin und Gesamtdrehimpuls}

\begin{definitionbox}
Der \emph{Spin} ist ein intrinsischer Drehimpuls eines Teilchens.
\end{definitionbox}

\begin{conceptbox}
Fermionen haben halbzahligen Spin:
\[
    s=\frac{1}{2},\frac{3}{2},\ldots
\]

Bosonen haben ganzzahligen Spin:
\[
    s=0,1,2,\ldots
\]
\end{conceptbox}

\begin{definitionbox}
Der \emph{Gesamtdrehimpuls} $J$ eines Teilchens oder Zustands setzt sich aus Spin
und Bahndrehimpuls zusammen.
\end{definitionbox}

\begin{conceptbox}
Drehimpuls ist bei allen fundamentalen Wechselwirkungen erhalten.
Bei Zerfaellen und Reaktionen muss daher der gesamte Drehimpuls vor und nach dem
Prozess zusammenpassen.
\end{conceptbox}

\section{Paritaet}

Paritaet beschreibt das Verhalten eines Zustands unter Raumspiegelung.

\begin{definitionbox}
Die \emph{Paritaetsoperation} spiegelt die Ortskoordinaten:
\[
    \vec r \rightarrow -\vec r.
\]
\end{definitionbox}

\begin{definitionbox}
Die \emph{Paritaet} $P$ eines Zustands ist der Eigenwert unter Raumspiegelung:
\[
    P=+1
\]
fuer gerade Paritaet und
\[
    P=-1
\]
fuer ungerade Paritaet.
\end{definitionbox}

\begin{conceptbox}
Bei einem Zweiteilchensystem haengt die Gesamtparitaet von den intrinsischen
Paritaeten der Teilchen und vom Bahndrehimpuls $L$ ab.
\end{conceptbox}

\begin{formulabox}
Fuer zwei Teilchen gilt:
\[
    P_{\mathrm{ges}}
    =
    P_1 P_2 (-1)^L.
\]
\end{formulabox}

\begin{notebox}
Fuer ein Fermion-Antifermion-System wie ein Meson kommt die intrinsische Paritaet
von Quark und Antiquark mit entgegengesetztem Vorzeichen ins Spiel.
Im einfachen Quarkmodell ergibt sich fuer Mesonen haeufig:
\[
    P=(-1)^{L+1}.
\]
\end{notebox}

\begin{examtipbox}
Paritaet wird von starker und elektromagnetischer Wechselwirkung erhalten.

Die schwache Wechselwirkung verletzt Paritaet.
Das ist ein wichtiges Kennzeichen schwacher Prozesse.
\end{examtipbox}

\section{C-Paritaet}

\begin{definitionbox}
Die \emph{Ladungskonjugation} $C$ vertauscht Teilchen und Antiteilchen.
\end{definitionbox}

\begin{conceptbox}
Nicht jeder Zustand hat eine definierte C-Paritaet.
Eine definierte C-Paritaet ist vor allem fuer neutrale Teilchen-Antiteilchen-Systeme
sinnvoll.
\end{conceptbox}

\begin{notebox}
Fuer Quarkonium-Zustaende kann C-Paritaet wichtig sein.
Quarkonia werden spaeter in einem eigenen Kapitel behandelt.
\end{notebox}

\begin{examtipbox}
C-Paritaet ist spezieller als elektrische Ladung oder Baryonenzahl.
Man sollte sie vor allem bei neutralen Teilchen-Antiteilchen-Systemen beachten.
\end{examtipbox}

\section{Erhaltungssaetze nach Wechselwirkung}

Die wichtigsten Erhaltungssaetze haengen von der beteiligten Wechselwirkung ab.

\begin{center}
\begin{tabular}{|p{4.2cm}|p{3.2cm}|p{3.2cm}|p{3.2cm}|}
\hline
Groesse & starke WW & EM-WW & schwache WW \\
\hline
Energie, Impuls & erhalten & erhalten & erhalten \\
\hline
elektrische Ladung $Q$ & erhalten & erhalten & erhalten \\
\hline
Baryonenzahl $B$ & erhalten & erhalten & erhalten \\
\hline
Leptonenzahl $L$ & erhalten & erhalten & erhalten \\
\hline
Strangeness $S$ & erhalten & erhalten & nicht immer erhalten \\
\hline
Quark-Flavor & erhalten & erhalten & kann sich aendern \\
\hline
Isospin & naeherungsweise erhalten & nicht exakt erhalten & nicht allgemein erhalten \\
\hline
Paritaet $P$ & erhalten & erhalten & verletzt \\
\hline
\end{tabular}
\end{center}

\begin{notebox}
Diese Tabelle ist eine Orientierung fuer typische Aufgaben.
Bei fortgeschrittener Teilchenphysik gibt es feinere Unterscheidungen, zum Beispiel
Neutrinooszillationen und CP-Verletzung.
Fuer diese Vorlesung ist die obige Einteilung die wichtigste Grundlage.
\end{notebox}

\section{Indikatoren fuer die beteiligte Wechselwirkung}

\begin{conceptbox}
Man kann die beteiligte Wechselwirkung oft aus Reaktionsmerkmalen erkennen.

Starke Wechselwirkung:
\begin{itemize}
    \item sehr schnelle Prozesse,
    \item Hadronen beteiligt,
    \item Strangeness wird bei Erzeugung erhalten,
    \item keine Leptonen als notwendige Endprodukte.
\end{itemize}

Elektromagnetische Wechselwirkung:
\begin{itemize}
    \item Photonen beteiligt,
    \item elektrische Ladung koppelt an Photon,
    \item Flavor bleibt erhalten,
    \item Paritaet bleibt erhalten.
\end{itemize}

Schwache Wechselwirkung:
\begin{itemize}
    \item Neutrinos koennen auftreten,
    \item Quark-Flavor kann sich aendern,
    \item Strangeness kann verletzt werden,
    \item Paritaet kann verletzt werden,
    \item Lebensdauern oft deutlich laenger als bei starken Zerfaellen.
\end{itemize}
\end{conceptbox}

\begin{examtipbox}
Wenn in einer Reaktion ein Neutrino auftritt, ist fast immer die schwache
Wechselwirkung beteiligt.

Wenn sich ein Quark-Flavor aendert, ist die schwache Wechselwirkung beteiligt.
\end{examtipbox}

\section{Beispiel: Neutronzerfall}

\begin{examplebox}
Betrachte:
\[
    n \rightarrow p + e^- + \overline{\nu}_e.
\]

Elektrische Ladung:
\[
    0 = 1 - 1 + 0.
\]

Baryonenzahl:
\[
    1 = 1 + 0 + 0.
\]

Leptonenzahl:
\[
    0 = 0 + 1 - 1.
\]

Auf Quarkebene:
\[
    d \rightarrow u + W^-.
\]

Der Prozess ist schwach, weil sich ein Quark-Flavor aendert.
\end{examplebox}

\section{Beispiel: Elektron-Positron-Annihilation}

\begin{examplebox}
Betrachte:
\[
    e^- + e^+ \rightarrow \gamma + \gamma.
\]

Elektrische Ladung:
\[
    -1+1=0.
\]

Leptonenzahl:
\[
    1-1=0.
\]

Die Photonen haben:
\[
    Q=0,
    \qquad
    L=0.
\]

Der Prozess ist elektromagnetisch.
\end{examplebox}

\section{Beispiel: Starke Erzeugung von Strange-Teilchen}

\begin{examplebox}
Betrachte:
\[
    \pi^- + p \rightarrow K^0 + \Lambda^0.
\]

Strangeness links:
\[
    S(\pi^-)+S(p)=0+0=0.
\]

Strangeness rechts:
\[
    S(K^0)+S(\Lambda^0)=+1-1=0.
\]

Strangeness ist erhalten.
Das passt zu starker Wechselwirkung.
\end{examplebox}

\begin{conceptbox}
Seltsame Teilchen koennen stark paarweise erzeugt werden.
Sie zerfallen aber haeufig schwach, weil einzelne Strangeness-Aenderungen in
Zerfaellen nur durch die schwache Wechselwirkung moeglich sind.
\end{conceptbox}

\section{Beispiel: Schwacher Strange-Zerfall}

\begin{examplebox}
Betrachte:
\[
    K^+ \rightarrow \mu^+ + \nu_\mu.
\]

Ladung:
\[
    +1 = +1 + 0.
\]

Leptonenzahl:
\[
    0 = -1 + 1.
\]

Strangeness:
\[
    S(K^+)=+1,
\]
aber
\[
    S(\mu^+)+S(\nu_\mu)=0.
\]

Strangeness ist also nicht erhalten.
Der Zerfall ist schwach.
\end{examplebox}

\section{Systematisches Vorgehen bei Reaktionen}

\begin{examtipbox}
Gehe bei Teilchenreaktionen systematisch vor:

\begin{enumerate}
    \item Schreibe alle Teilchen mit ihren Quantenzahlen auf.
    \item Pruefe elektrische Ladung.
    \item Pruefe Baryonenzahl.
    \item Pruefe Leptonenzahl.
    \item Pruefe Strangeness und andere Flavor-Quantenzahlen.
    \item Pruefe, ob Paritaet relevant ist.
    \item Entscheide, welche Wechselwirkung den Prozess erlauben kann.
\end{enumerate}
\end{examtipbox}

\begin{conceptbox}
Wenn alle Flavor-Quantenzahlen erhalten sind und nur Hadronen beteiligt sind, kann
der Prozess stark sein.

Wenn ein Photon beteiligt ist und Flavor erhalten bleibt, kann der Prozess
elektromagnetisch sein.

Wenn Flavor verletzt wird oder Neutrinos beteiligt sind, ist der Prozess schwach.
\end{conceptbox}

\section{Mini-Zusammenfassung}

\begin{notebox}
Die wichtigsten Punkte dieses Kapitels sind:
\begin{itemize}
    \item Quantenzahlen beschreiben Eigenschaften von Teilchen.
    \item Elektrische Ladung ist in allen fundamentalen Prozessen erhalten.
    \item Quarks haben Baryonenzahl
    \[
        B=\frac{1}{3}.
    \]
    \item Baryonen haben
    \[
        B=1.
    \]
    \item Mesonen haben
    \[
        B=0.
    \]
    \item Leptonen haben Leptonenzahl
    \[
        L=+1.
    \]
    \item Antileptonen haben
    \[
        L=-1.
    \]
    \item Strangeness des strange-Quarks ist
    \[
        S(s)=-1.
    \]
    \item Strangeness des strange-Antiquarks ist
    \[
        S(\overline{s})=+1.
    \]
    \item Starke und elektromagnetische Wechselwirkung erhalten Flavor.
    \item Schwache Wechselwirkung kann Flavor aendern.
    \item Hyperladung ist im einfachen strange-Hadron-Kontext
    \[
        Y=B+S.
    \]
    \item Die Gell-Mann--Nishijima-Formel lautet:
    \[
        Q=I_3+\frac{Y}{2}.
    \]
    \item Paritaet wird von starker und elektromagnetischer Wechselwirkung erhalten.
    \item Schwache Wechselwirkung verletzt Paritaet.
\end{itemize}
\end{notebox}

\section{Pruefungscheck}

\begin{examtipbox}
Nach diesem Kapitel solltest du folgende Fragen beantworten koennen:
\begin{enumerate}
    \item Was ist eine Quantenzahl?
    \item Was ist der Unterschied zwischen additiven und multiplikativen Quantenzahlen?
    \item Warum ist elektrische Ladung ein wichtiger erster Check?
    \item Welche Baryonenzahl hat ein Quark?
    \item Welche Baryonenzahl hat ein Antiquark?
    \item Welche Baryonenzahl haben Baryonen, Antibaryonen und Mesonen?
    \item Was ist Leptonenzahl?
    \item Warum tritt beim Beta-Minus-Zerfall ein Antineutrino auf?
    \item Was bedeutet Flavor?
    \item Was ist Strangeness?
    \item Welches Vorzeichen hat die Strangeness des strange-Quarks?
    \item Wann ist Strangeness erhalten?
    \item Wann kann Strangeness verletzt werden?
    \item Was ist Isospin?
    \item Was ist die dritte Isospinkomponente $I_3$?
    \item Was ist Hyperladung?
    \item Wie lautet die Gell-Mann--Nishijima-Formel?
    \item Was ist Paritaet?
    \item Welche Wechselwirkungen erhalten Paritaet?
    \item Welche Wechselwirkung verletzt Paritaet?
    \item Woran erkennt man starke, elektromagnetische und schwache Prozesse?
    \item Wie prueft man systematisch, ob eine Teilchenreaktion erlaubt ist?
\end{enumerate}
\end{examtipbox}